\section{Label Overlap Stress Test -- 1D Arrays and DPTables}

\subsection{Scenario A: Three annotations on the same cell with long labels}

\begin{animation}[id="overlap-same-cell", label="3 annotations on one cell"]
\shape{dp}{DPTable}{n=6, data=[0,7,2,12,8,15], labels="0..5", label="$dp$"}

\step
\narrate{Three annotations all targeting dp[3] from different sources with long labels.}
\annotate{dp.cell[3]}{label="dp[0]+cost(0,3)=42", arrow_from="dp.cell[0]", color=info}
\annotate{dp.cell[3]}{label="dp[1]+cost(1,3)=35", arrow_from="dp.cell[1]", color=good}
\annotate{dp.cell[3]}{label="dp[2]+cost(2,3)=28", arrow_from="dp.cell[2]", color=warn}

\step
\narrate{Now add a fourth annotation from the right side to make it worse.}
\annotate{dp.cell[3]}{label="dp[5]+cost(5,3)=99", arrow_from="dp.cell[5]", color=error}
\end{animation}

\subsection{Scenario B: Adjacent cells each with their own annotations}

\begin{animation}[id="overlap-adjacent", label="Adjacent cell annotations"]
\shape{dp2}{DPTable}{n=6, data=[0,3,5,8,12,17], labels="0..5", label="$dp$"}

\step
\narrate{Every adjacent pair has an annotation -- labels crowd the space between cells.}
\annotate{dp2.cell[1]}{label="dp[0]+3=3", arrow_from="dp2.cell[0]", color=info}
\annotate{dp2.cell[2]}{label="dp[1]+2=5", arrow_from="dp2.cell[1]", color=info}
\annotate{dp2.cell[3]}{label="dp[2]+3=8", arrow_from="dp2.cell[2]", color=info}
\annotate{dp2.cell[4]}{label="dp[3]+4=12", arrow_from="dp2.cell[3]", color=info}
\annotate{dp2.cell[5]}{label="dp[4]+5=17", arrow_from="dp2.cell[4]", color=info}

\step
\narrate{Now add reverse-direction annotations that cross the forward ones.}
\annotate{dp2.cell[1]}{label="dp[3]-cost=5", arrow_from="dp2.cell[3]", color=warn}
\annotate{dp2.cell[2]}{label="dp[4]-cost=7", arrow_from="dp2.cell[4]", color=warn}
\annotate{dp2.cell[3]}{label="dp[5]-cost=9", arrow_from="dp2.cell[5]", color=warn}
\end{animation}

\subsection{Scenario C: Very short array (size=3) with full-span annotations}

\begin{animation}[id="overlap-short-array", label="Short array full-span annotations"]
\shape{s}{DPTable}{n=3, data=[0,5,9], labels="0..2", label="$dp$"}

\step
\narrate{Tiny array with wide-spanning annotations -- labels have nowhere to go.}
\annotate{s.cell[2]}{label="dp[0]+cost(0,2)=9", arrow_from="s.cell[0]", color=good}
\annotate{s.cell[2]}{label="dp[1]+cost(1,2)=4", arrow_from="s.cell[1]", color=info}
\annotate{s.cell[1]}{label="dp[0]+cost(0,1)=5", arrow_from="s.cell[0]", color=warn}
\end{animation}

\subsection{Scenario D: Very long array (size=12) with 10+ cell spans}

\begin{animation}[id="overlap-long-array", label="Long array with long-span annotations"]
\shape{big}{DPTable}{n=12, data=[0,3,5,8,12,17,21,26,30,35,40,48], labels="0..11", label="$dp$"}

\step
\narrate{Long-span annotations across 10+ cells on a 12-cell array.}
\annotate{big.cell[11]}{label="dp[0]+cost(0,11)=48", arrow_from="big.cell[0]", color=good}
\annotate{big.cell[10]}{label="dp[1]+cost(1,10)=37", arrow_from="big.cell[1]", color=info}
\annotate{big.cell[11]}{label="dp[2]+cost(2,11)=43", arrow_from="big.cell[2]", color=warn}

\step
\narrate{Add more crossing annotations at different heights.}
\annotate{big.cell[6]}{label="dp[0]+cost(0,6)=21", arrow_from="big.cell[0]", color=error}
\annotate{big.cell[8]}{label="dp[3]+cost(3,8)=22", arrow_from="big.cell[3]", color=path}
\annotate{big.cell[5]}{label="dp[10]+cost(10,5)=45", arrow_from="big.cell[10]", color=muted}
\end{animation}
